% =========================================================================
% When Does Feedback Help? Planning and Model Mismatch in Hessian-Free
% Newton Control
%
% **표현 계층이다.** 내용의 source of truth 는 `paper/draft.md` 다.
% 주장 강도는 `paper/claim_ledger.md` 가 통제하고, 인용 내용 검증은
% `paper/CITATIONS.md` 에 있다.
%
% 규칙
%   1. 표의 숫자를 이 파일이나 sections/ 에 쓰지 않는다.
%      `scripts/make_tables.py` 가 `tables/*.tex` 를 생성한다
%   2. 그림은 `scripts/make_figures.py` 가 생성한다
%   3. 새 주장을 여기서 만들지 않는다. draft.md 에 없는 문장을 넣지 않는다
%
% 빌드
%   pdflatex main && bibtex main && pdflatex main && pdflatex main
% =========================================================================
\documentclass[11pt]{article}

\usepackage[utf8]{inputenc}
\usepackage[T1]{fontenc}
\usepackage[english]{babel}
\usepackage{lmodern}
\usepackage{microtype}

\usepackage[margin=1in]{geometry}
\usepackage{amsmath,amssymb}
\usepackage{booktabs}
\usepackage{graphicx}
\usepackage{subcaption}
\usepackage{xcolor}
\usepackage[numbers,sort&compress]{natbib}
\usepackage[hidelinks]{hyperref}

\graphicspath{{figures/}}

% 컨트롤러 이름을 본문에서 일관되게 쓴다.
\newcommand{\ctrl}[1]{\texttt{#1}}
\newcommand{\GE}{\text{GE}}
\newcommand{\JE}{J_E}

% 사전 등록 게이트 표기.
\newcommand{\gate}[1]{\textsf{#1}}

\title{When Does Feedback Help?\\
Planning and Model Mismatch in Hessian-Free Newton Control}

\author{%
  % 저자 정보는 제출 직전에 채운다.
  Anonymous
}
\date{}

\begin{document}
\maketitle

\input{sections/00_abstract}
\input{sections/01_introduction}
\input{sections/02_setup}
\input{sections/03_controllers}
\input{sections/04_protocol}
\input{sections/05_eligibility}
\input{sections/06_selection}
\input{sections/07_heldout}
\input{sections/08_decomposition}
\input{sections/09_conditioning}
\input{sections/10_mismatch}
\input{sections/11_ablation}
\input{sections/12_nopolicy}
\input{sections/13_discussion}
\input{sections/14_limitations}
\section{Reproducibility}
\label{sec:reproducibility}

\subsection{Execution environment}

\begin{quote}
All Stage~2 optimization and planning experiments were executed on CPU. The
GPU-derived cost-model files in the repository originate from an earlier Stage~1
calibration and were not used in Stage~2; Stage~2 GE accounting used HVP-equivalent
oracle counts.
\end{quote}

\begin{quote}
Search GE measures simulated oracle work rather than wall-clock GPU compute.
\end{quote}

Execution was pinned to a single thread. Wall-clock time is used in no verdict.

\subsection{Artifacts}

\begin{center}
\begin{tabular}{ll}
\toprule
Artifact & Location \\
\midrule
Reproduction commands       & \texttt{docs/reproduce.md} \\
Generated result tables     & \texttt{docs/results\_stage2.md} \\
Raw checksums (SHA-256)     & \texttt{paper/evidence\_map.md} \\
Claim ledger                & \texttt{paper/claim\_ledger.md} \\
Citation content checklist  & \texttt{paper/CITATIONS.md} \\
Protocol (decisions D1--D32) & \texttt{docs/experiment\_protocol.md} \\
Tag                         & \texttt{protocol-freeze-stage2-v1} \\
Tests                       & 455 \\
\bottomrule
\end{tabular}
\end{center}

All numbers in this manuscript originate in \texttt{docs/results\_stage2.md}, whose
sources are pinned by the SHA-256 values in \texttt{paper/evidence\_map.md}. The
tables in \texttt{paper/tables/} are generated by
\texttt{scripts/make\_tables.py} from the same raw records and the same median
convention. \textbf{Numbers are not edited by hand in this manuscript.}

The result store skips completed runs. Because identity is separated into three
layers, editing aggregation code or documentation does not cause the optimizer to
rerun.

\subsection*{Privacy note}

\begin{quote}
Hostnames in committed public artifacts were replaced with stable pseudonymous labels
before public release. This sanitization did not modify experimental configurations,
results, checksums of raw numerical records, or scientific conclusions.
\end{quote}

Git ancestry still contains the earlier hostname, so this release is not described as
fully anonymized in history.


\bibliographystyle{plainnat}
\bibliography{references}

\end{document}
